% La comparaison situe la bibliothèque et ses limites.

% MIN 25:30 à 26:30. Comparaison des bibliothèques retenues pour le rapport.
\begin{frame}{Face à cinq bibliothèques comparables}
    Pour la comparaison, j'ai retenu ByzFL, FedLab, Blades,
    FL-Byzantine-Library et ByzPy, dont les composants d'agrégation byzantine
    sont ouverts et comparables.

    \medskip
    {\renewcommand{\arraystretch}{1.2}
    \begin{tabular}{@{}p{0.18\textwidth}p{0.34\textwidth}p{0.40\textwidth}@{}}
        \textbf{Critère} & \textbf{Krum} & \textbf{Dans la comparaison} \\
        \hline
        Catalogue & 11 règles et 5 attaques. & Certaines bibliothèques proposent plus de composants. \\
        Zéro-copie & \code{Model} partage la mémoire. & L'accès copie les valeurs ou n'a pas d'équivalent. \\
        Protocoles & Trois protocoles documentés. & Cadres génériques ou suites. \\
        Relance & Recalcul des résultats périmés. & ByzFL saute les exécutions terminées.
    \end{tabular}}

    \medskip
    J'ai analysé ces composants.\\
    {\small Comptes d'août 2026, audités en septembre 2026.}
\end{frame}

% MIN 26:30 à 27:00. Limites du catalogue et des conditions d'exécution.
\begin{frame}{Les limites actuelles}
    \begin{itemize}
        \item Le catalogue reste ciblé. Certaines bibliothèques proposent plus de composants.
        \item Les pré-agrégateurs ne sont pas inclus. Appliqués avant l'agrégation,
              ils peuvent borner les normes ou réduire le bruit.
        \item Les expériences s'exécutent séquentiellement sur une seule machine.
              Les échanges entre participants sont simulés. On évalue donc les
              protocoles, pas un déploiement distribué réel.
    \end{itemize}
\end{frame}
